\section{Sharp comparison for the standard spectral barriers}
\label{sec:sharp-centrality}
The reduction in \cref{thm:water-filling} applies to all products considered
there. An active channel means a positive singular value of a matrix
objective, or the absolute value of a nonzero Jordan eigenvalue. Denote
these weights by $w_i>0$ and their factor scales by $\alpha_i>0$,
$i=1,\ldots,r$. Signs and spectral frames are fixed by the objective;
zero channels remain zero. The metric assertions in this section require
only positive scales. Statements about standard self-concordance and
bounded local moves additionally require $\alpha_i\geq1$.

\subsection{The scalar velocity and the exact prefix constant}
Put
\begin{equation}
 q(z)=\frac{z}{1+\sqrt{1+z^2}},\qquad
 H(s)=\rho(q(e^s)),\qquad
 v(s)=H'(s)=\sqrt{1-\frac1{\sqrt{1+e^{2s}}}}.
 \label{eq:standard-profile}
\end{equation}
Here $H$ is a strictly increasing bijection from $\R$ to $(0,\infty)$.
The central coordinates, in the objective frame, are
\begin{equation}
 x_i(\eta)=q(\eta w_i/\alpha_i),\quad
 a_i=\log(\alpha_i/w_i),\quad
 y_i(s)=\rho(x_i(e^s))=H(s-a_i).
 \label{eq:central-coordinates}
\end{equation}
For a matrix factor of scale $\alpha$, this is equivalently
\[
 X(\eta)=\frac\eta\alpha
 \left(I+\sqrt{I+(\eta/\alpha)^2CC^*}\right)^{-1}C.
\]
These formulas follow by minimizing each scalar function
$\alpha_i b(x)-\eta w_ix$: its derivative equation is
$2\alpha_i x/(1-x^2)=\eta w_i$. The spectral trace inequalities used in
\cref{thm:water-filling} show that the aligned point minimizes in the
whole domain. Strict convexity proves uniqueness, including tied and
zero objective weights. Differentiation gives
$dx/ds=x(1-x^2)/(1+x^2)$ and then \eqref{eq:standard-profile}.
In particular, $v$ increases strictly from zero to one and
\begin{equation}
 v(s)\leq e^s/\sqrt2\ (s\leq0),\qquad
 1-v(s)\leq e^{-s}\ (s\geq0),\qquad
 H(s)=s+\kappa+o(1)\ (s\to\infty),
 \label{eq:profile-tails}
\end{equation}
for a finite constant $\kappa$. The last assertion follows by integrating
the first two bounds. They also show $H(s)\leq1/\sqrt2+\max(s,0)$.

Order channels by nondecreasing $a_i$, breaking ties arbitrarily, and set
\begin{equation}
 S_0=0,\quad S_i=\sum_{j\leq i}\alpha_j,\quad
 d_i=\sqrt{S_i}-\sqrt{S_{i-1}},\quad
 \Gamma_{\boldsymbol\alpha,a}^2=\sum_{i=1}^r\frac{d_i^2}{\alpha_i}.
 \label{eq:weighted-Gamma}
\end{equation}
The value of this coefficient can depend on the ordering chosen inside a
tied group; any such ordering gives the following bound. Its sharpness
statement allows the objective to choose strictly separated thresholds.
For equal scales we write
\begin{equation}
 \Gamma_r^2=\sum_{i=1}^r(\sqrt i-\sqrt{i-1})^2
 =\tfrac14\log r+O(1).
 \label{eq:Gamma-r}
\end{equation}

\begin{theorem}[Exact worst comparison for a prescribed scale order]
\label{thm:sharp-prefix}
For every central subarc, including the analytic-center limit,
\begin{equation}
 d_\Phi(x(\eta_0),x(\eta_1))
 \leq L_\CP(\eta_0,\eta_1)
 \leq\Gamma_{\boldsymbol\alpha,a}
           d_\Phi(x(\eta_0),x(\eta_1)),\qquad 0\leq\eta_0<\eta_1.
 \label{eq:sharp-prefix}
\end{equation}
For each fixed ordered positive scale vector, the supremum of the
analytic-center-to-central-endpoint ratio, over objectives enforcing this
strict activation order and over terminal parameters, is exactly
$\Gamma_{\boldsymbol\alpha,a}$. It is already realized as a supremum on
a product of scalar intervals.
\end{theorem}
\begin{proof}
For $h_1\geq\cdots\geq h_r\geq0$, put $h_{r+1}=0$. Decomposition into
initial-coordinate indicator vectors and the triangle inequality give
\begin{equation}
 \left(\sum_i\alpha_ih_i^2\right)^{1/2}
 \leq\sum_j\sqrt{S_j}(h_j-h_{j+1})=\sum_i d_ih_i.
 \label{eq:prefix-inequality}
\end{equation}
Apply this to $h_i=v(s-a_i)$ and integrate over the parameter interval.
Weighted Cauchy--Schwarz bounds the result by
$\Gamma_{\boldsymbol\alpha,a}(\sum_i\alpha_i(\Delta y_i)^2)^{1/2}$.
The latter norm is the exact endpoint distance by
\cref{thm:exact-distance}. In a Jordan factor, the fixed signs and order
of absolute objective eigenvalues ensure a common ordered signed frame
throughout the path; in a matrix factor the singular frame is likewise
coordered. Thus this equality does apply to central subarcs.

For sharpness set
$c_i=d_i/\alpha_i=1/(\sqrt{S_i}+\sqrt{S_{i-1}})$, a strictly decreasing
sequence. For $M\to\infty$, choose $u_i=H^{-1}(Mc_i)$,
$w_i=\alpha_i e^{u_i}$, and terminal $s=0$. The endpoint distance is
exactly $M\Gamma_{\boldsymbol\alpha,a}$. Replacing each
$v(s+u_i)$ by $\mathbf1_{s>-u_i}$ changes the integrated weighted norm
on $(-\infty,0]$ by at most
$\sum_i\sqrt{\alpha_i}\int_\R|v-\mathbf1_{s>0}|\dd s$, independent
of $M$. For large $M$, all $u_i$ are positive, and the step-function
length is $\sum_i d_iu_i$. By \eqref{eq:profile-tails},
$u_i=Mc_i-\kappa+o(1)$, so this length is
$M\Gamma_{\boldsymbol\alpha,a}^2+O_{\boldsymbol\alpha}(1)$.
Taking the ratio proves the supremum. Finally,
$(\sqrt i-\sqrt{i-1})^2=1/(4i)+O(i^{-2})$ proves
\eqref{eq:Gamma-r}.
\end{proof}

The decreasing rearrangement norm in \eqref{eq:prefix-inequality} is a
finite weighted form of a classical Lorentz norm~\cite{Lorentz1951}.
The elementary prefix inequality itself is not a novelty claim; the
sharp central-path realization is the point of \cref{thm:sharp-prefix}.
Unlike the general comparison \eqref{eq:NN-comparison}, its equal-scale
constant depends on active objective rank rather than ambient dimension.
The sharp limiting objectives can have large dynamic range. A separate
rational-input construction is given in \cref{sec:discrete}.

\begin{proposition}[Scale range and activation order]
\label{prop:scale-order}
For $r\geq2$,
\begin{equation}
 \Gamma_{\boldsymbol\alpha,a}^2
 \leq1+(r-1)\tanh\left(\frac{\log(S_r/S_1)}{4(r-1)}\right)
 \leq\min\{r,\,1+\tfrac14\log(S_r/S_1)\}.
 \label{eq:scale-frontier}
\end{equation}
Equality in the first inequality holds precisely when the cumulative
ratios $S_i/S_{i-1}$, $i\geq2$, are equal. Among orders of a fixed
multiset of scales, nondecreasing scales maximize $\Gamma$ and
nonincreasing scales minimize it.
\end{proposition}
\begin{proof}
The first term of \eqref{eq:weighted-Gamma} is one; its $i$th term,
$i\geq2$, is $\tanh(\frac14\log(S_i/S_{i-1}))$.
Strict concavity of $\tanh$ on $(0,\infty)$ gives the first inequality
and its equality condition. The bounds $\tanh z\leq\min(1,z)$ give
the second.
For the ordering assertion, interchange adjacent scales $a\leq b$ with
preceding mass $S>0$. The two logarithmic increments have common sum
$T=\log((S+a+b)/S)$. Their first increments are
$u=\log((S+a)/S)\leq v=\log((S+b)/S)$, and $u+v\geq T$ because
$(S+a)(S+b)\geq S(S+a+b)$. Thus
$|u-T/2|\leq|v-T/2|$. Concavity makes the symmetric sum
$\tanh(u/4)+\tanh((T-u)/4)$ at least its value at $v$.
If $S=0$, the first contribution is one and the second is larger with
the smaller scale first. Adjacent exchanges prove both assertions.
\end{proof}
For example, cumulative-geometric masses with a fixed ratio greater than
one give $\Gamma=\Theta(\sqrt r)$. Conversely, this order requires
$\log(S_r/\alpha_{\min})=\Omega(r)$. With $\alpha_i\geq1$ and
polynomially bounded total scale, the overhead remains $O(\sqrt{\log r})$.
For $r=1$, $\Gamma=1$ without a scale restriction.

\subsection{The exact scalar comparison with the closest accurate endpoint}
Define $p(y)=b'(\rho^{-1}(y))$ for $y\geq0$, using the functional inverse
of $\rho$. Then $p'(y)=\rho'(\rho^{-1}(y))>0$ and $p$ is strictly convex
on $(0,\infty)$.
\begin{theorem}[Uniform scalar dilation and the accuracy comparison]
\label{thm:scalar-dilation}
The exact least constant satisfying $p(cy)\geq yp'(y)$ for all $y>0$ is
\begin{equation}
 c_\star=\max_{y>0}\frac{p^{-1}(yp'(y))}{y},\qquad
 \frac{68743}{50000}<c_\star<\frac{69}{50}.
 \label{eq:cstar}
\end{equation}
The maximum is attained at a finite positive argument. For every product
under consideration and $0<\eps<\sum_iw_i$,
\begin{equation}
 L_\CP(\eps)\leq c_\star\Gamma_{\boldsymbol\alpha,a}L_\opt(\eps).
 \label{eq:same-accuracy}
\end{equation}
The constant $c_\star$ is sharp for this scalar inequality; no assertion
is made that its product with $\Gamma$ is the sharp arclength ratio.
\end{theorem}
\begin{proof}
Use $v=\sqrt2x/\sqrt{1+x^2}\in(0,1)$ and write
\begin{align}
 Y(v)&=2\operatorname{artanh}v-
             \sqrt2\operatorname{artanh}(v/\sqrt2)=\rho(x),\notag\\
 A(v)&=\frac{\sqrt{2-v^2}}{1-v^2}=p'(Y(v)),\qquad
 P(v)=vA(v)=p(Y(v)),\qquad
 Y'(v)=\frac2{(1-v^2)(2-v^2)}.
 \label{eq:scalar-elementary}
\end{align}
These identities follow by differentiation and equality at zero. Put
\[
 W(v)=\sqrt{1-(1+Y(v)^2A(v)^2)^{-1/2}}.
\]
The ratio defining $c_\star$ is $Y(W(v))/Y(v)$. As $v\downarrow0$, $Y(v)=v+O(v^3)$ and this ratio
tends to one. As $v\uparrow1$, $Y(v)=\log(1/(1-v))+O(1)$ and
$1-W(v)=\Theta((1-v)/Y(v))$, so
$Y(W(v))=Y(v)+\log Y(v)+O(1)$ and the ratio again tends to one.
Strict convexity gives $yp'(y)>p(y)$, so the ratio exceeds one inside
the interval. Continuity and these endpoint limits prove attainment.
\Cref{app:scalar-certificate} proves both strict rational bounds.

Let $y_i^\star=\rho(x_i^\star)$ be the exact allocation in
\cref{thm:water-filling}, with multiplier $\lambda>0$.
Its equations are $2\alpha_i y_i^\star p'(y_i^\star)=\lambda w_i$.
At central parameter $\eta=\lambda/2$, the central coordinates obey
$\alpha_i p(y_i^{\mathrm c})=\lambda w_i/2$, hence
\[
 p(y_i^{\mathrm c})=y_i^\star p'(y_i^\star),\qquad
 y_i^\star\leq y_i^{\mathrm c}\leq c_\star y_i^\star.
\]
This center is accurate. The first accurate center occurs no later and
is no larger in any active coordinate. Its distance from zero is at most
$c_\star L_\opt(\eps)$. Apply \cref{thm:sharp-prefix}.
\end{proof}
The formula \eqref{eq:scalar-elementary} gives a direct one-variable way
to evaluate the constant. Numerical evaluation suggests
$c_\star\approx1.37486420044$; the theorem uses only its variational
definition and the rational certificates. Neither uniqueness of its
maximizer nor a decimal accuracy guarantee is needed.
